\subsection{Fitting and timing for the earlier comparisons}
\label{app:asrr-tuning-details}

Unaugmented Omniglot uses disjoint sets of 1028 source, 172 development and 423
target characters. We use 784--16--5 and 784--64--5 ReLU networks. Each target
episode samples five characters and a rotation in multiples of 90
degrees for each character. Its 20 drawings are divided into 10 fitting,
five numerical-selection and five query images. The banks use the 16
source tasks with lowest numerical-validation cross-entropy among 128
earlier five-way fits. Each source class has five fitting, five
numerical-selection and ten structure-selection drawings. The final
comparison at both widths uses the same 1024 query episodes as the initial
WSR run.

On HAR, the 561--16--6 ReLU network receives the original features.
We retain every second window within each subject and activity before
drawing episodes. Each source task has 366 fitting, 78 numerical-selection
and 78 structure-selection windows. The 96 development episodes select
the numerical rule for each frozen or eligible 128-coordinate bank.
We screen learning rates 0.001, 0.003, 0.01 and 0.03 with activation
quantiles 0.1, 0.5 and 0.9 and up to 1000 Adam steps, then compare
ensembles of 1, 4 and 16 predictions for the retained settings.
The four source tasks each pool four participants. Five other
participants provide development episodes. The nine held-out
participants provide 180 test episodes with three fitting and five
query windows per activity. Table~\ref{tab:recovery-har} reports
accuracy, cross-entropy and source construction times.
\begin{table}[htbp]
\centering\small
\caption{UCI HAR on 180 common episodes. Accuracy gives 95\% intervals over the nine subject means.}
\label{tab:recovery-har}
\begin{tabular}{lrrrr}
\toprule
Method & Budget & Accuracy (\%) & CE & Source time (s) \\
\midrule
Dense & 8,976 & $86.43\pm3.10$ & 0.3753 & \textemdash \\
aSRR & 128 & $91.50\pm2.79$ & 0.2665 & 7.08 \\
aSRR (shuffled groups) & 128 & $86.62\pm2.97$ & 0.3894 & \textemdash \\
LUT-Q & 128 & $88.70\pm3.22$ & 0.3412 & 87.96 \\
Yeh (MAP) & 128 & $88.93\pm2.93$ & 0.3376 & 569.2 \\
Yeh (soft) & 128 & $90.87\pm3.88$ & 0.4106 & 569.2 \\
MSR-FC & 128 & $89.43\pm3.05$ & 0.3358 & 150.9 \\
WSR & 128 & $85.02\pm4.58$ & 1.1387 & 6.82 \\
\bottomrule
\end{tabular}
\end{table}


LUT-Q, Yeh and MSR-FC use fully connected adaptations with the same
source observations. The respective source budgets are 2500 LUT-Q
updates, 400 Yeh outer updates with a 20-term Neumann approximation,
and 2000 MSR updates. On Omniglot up to eight candidates per native
method enter development, ranked first on source numerical validation.
LUT-Q and Yeh use signed dictionaries with a fixed zero. The continuous
variants transfer only the learned basis.

WSR starts from each source task's Dense fit and runs 500 full-batch
Adam steps at learning rate 0.003. At each step we apply the WSR proximal
map to the sorted first-layer weights using isotonic regression. Bisection
sets the strength for the prescribed group count. We apply rewinding
with coefficient 0.98~\cite{shakerinava2024weight} and retain the best
numerical-validation checkpoint from steps 0, 25, \ldots, 500. Each
dataset's development episodes select one of the source groupings when
several are available.

Source times come from construction logs and replays of fixed
configurations. aSRR and WSR start from saved Dense fits. The aSRR
timing in the sparse ReLU experiment covers both candidates. Yeh MAP
and soft share one source run. Archived native times include profiling
overhead. All times end at candidate generation, before selection and
transfer.

For the 64-hidden-unit Omniglot model, we fit the same 16 source tasks
with 12 Dense starts of 2500 updates. aSRR clusters the input profiles
into 64 clusters and assigns a separate code to each cluster in each
hidden neuron. WSR selects its grouping on 96 development episodes.
One of its 8016 proximal updates skipped the prescribed group count.
We resolved the tie with perturbations at float32 precision and used
adjacent representable values to separate five rounded code levels.
Yeh (soft) uses a single configuration, fixed by the development results
at width 16. It initializes from scratch, uses outer and inner learning
rates of 0.03 and 0.01, and sets both regularization coefficients to zero.
We run 400 outer updates with 250 initial and 20 subsequent inner steps.

All target fits process batches of 24 episodes. We time fresh numerical
fitting and within-task checkpoint selection on the first 96 test episodes
on one Tesla V100 PCIe (32 GiB). Peak host memory is process RSS, and
peak GPU memory is PyTorch's maximum allocated memory.

\subsection{Input clustering on wider Omniglot}
\label{app:omniglot-clustering}

Each algorithm partitions the same 784 input profiles from 16 source
tasks into 64 groups. We reuse the 4096-code budget, numerical fitting
rule and 1024 test episodes of Table~\ref{tab:omniglot-width64}.

Lloyd k-means~\cite{lloyd1982quantization} uses four starts of up to 25
updates. The authors' FP32 PyTorch e-means~\cite{bessa2026emeans} runs on
CPU with four starts each for \texttt{auto} and \texttt{quality} at five
epochs and \texttt{quality} at 25. We refill empty clusters from
high-residual profiles and retain the partition with the smallest
within-cluster sum of squares on source profiles.

e-means leads Lloyd by $0.09$ points (unadjusted paired 95\% CI
$[-0.21,0.39]$). Single linkage puts 689 of 784 inputs in one cluster.

\begin{table}[!htbp]
\centering\small
\caption{Input clustering at 4096 codes on 1024 shared Omniglot episodes.
Accuracy gives marginal 95\% intervals. Paired differences from Ward use
Bonferroni-adjusted intervals for five comparisons. Source times include
all runs from saved Dense fits. On 96 shared episodes, target fitting
takes 33.42--35.35 seconds with 0.44 GiB of peak allocated GPU memory.}
\label{tab:omniglot-clustering}
\begin{tabular}{lrrrr}
\toprule
Clusterer & Accuracy (\%) & $\Delta$ vs. Ward (pp) & 95\% CI (pp) & Source time (s) \\
\midrule
Ward & $83.55\pm0.54$ & \textemdash & \textemdash & 0.58 \\
Lloyd k-means & $83.88\pm0.53$ & $+0.32$ & $[-0.10,0.75]$ & 0.66 \\
e-means & $83.96\pm0.52$ & $+0.41$ & $[0.02,0.81]$ & 9.21 \\
Complete linkage & $83.38\pm0.53$ & $-0.18$ & $[-0.56,0.21]$ & 0.98 \\
Average linkage & $80.39\pm0.58$ & $-3.16$ & $[-3.69,-2.64]$ & 1.01 \\
Single linkage & $71.09\pm0.67$ & $-12.46$ & $[-13.26,-11.66]$ & 1.23 \\
\bottomrule
\end{tabular}
\end{table}
